\title{Pixelated Butterfly: Simple and Efficient Sparse Training for Neural Network Models}

\begin{abstract}
Overparameterized neural networks generalize well but are expensive to train. Ideally, one would like to reduce their computational cost while retaining their generalization benefits. Sparse model training is a simple and promising approach to achieve this, but there remain challenges as existing methods struggle with accuracy loss, slow training runtime, or difficulty in sparsifying all model components. The core problem is that searching for a sparsity mask over a discrete set of sparse matrices is difficult and expensive. To address this, our main insight is to optimize over a continuous superset of sparse matrices with a fixed structure known as products of butterfly matrices. As butterfly matrices are not hardware efficient, we propose simple variants of butterfly (block and flat) to take advantage of modern hardware. Our method (Pixelated Butterfly) uses a simple fixed sparsity pattern based on flat block butterfly and low-rank matrices to sparsify most network layers (e.g., attention, MLP). We empirically validate that Pixelated Butterfly is $3\times$ faster than butterfly and speeds up training to achieve favorable accuracy--efficiency tradeoffs. On the ImageNet classification and WikiText-103 language modeling tasks, our sparse models train up to 2.5$\times$ faster than the dense MLP-Mixer, Vision Transformer, and GPT-2 medium with no drop in accuracy.
\end{abstract}

\section{Introduction}
\label{sec:intro}

Recent results suggest that overparameterized neural networks generalize well ~\citep{belkin2019reconciling}, but they are expensive to train~\citep{kaplan2020scaling}. An ideal model should use less compute and memory while retaining the generalization benefits of large models. The simplest and most popular direction is to sparsify these models. This idea has a long history in machine learning~\citep{lecun1990optimal} and has driven fundamental progress in other fields such as statistics~\citep{tibshirani1996regression}, neuroscience~\citep{foldiak2003sparse}, and signal processing~\citep{candes2006stable}. However, despite significant efforts, {\em speeding up sparse training in wall-clock time without degrading accuracy} remains an unresolved problem.

While sparse training is an active research area, it has not seen wide adoption. First, it is difficult and expensive to find the sparsity pattern (the possible locations of the nonzeros) that could maintain the same level of accuracy of dense models. Many methods (pruning~\citep{lee2018snip}, lottery tickets~\citep{frankle2018lottery}, hashing~\citep{kitaev2020reformer}) maintain dynamic sparsity masks. However, the large overhead of evolving the sparsity mask can significantly slow down training and complicate the implementation. Indeed, these methods either require long cycles of pruning and retraining~\citep{frankle2018lottery}\footnote{State-of-the-art sparse training methods require up to 5$\times$ more training epochs compared to dense models~\citep{evci2020rigging}} or maintain expensive hash tables~\citep{chen2019slide}. Second, most existing methods adopt unstructured sparsity, which may be efficient in theory, but do not take into account the efficiency of training hardware such as GPUs (optimized for dense computation)\footnote{An unstructured sparse model with 1\% nonzero weights can be as slow as a dense model~\citep{hooker2020hardware}}. Finally, most methods target a single type of operation such as attention~\citep{child2019generating,zaheer2020bigbird}, whereas neural network (NN) models often compose different modules (attention, MLP), and in many applications the MLP layers are the main training bottleneck~\citep{wu2020lite}.

A better sparse training method should (i) be simple yet accurate, ideally with a static sparsity pattern, (ii) be fast by aligning sparsity pattern with available hardware, and (iii) have wide coverage of operators that applies to most NN layers. There are three technical challenges. First, we show that given a budget (e.g., total non-zeros in a matrix), it is NP-hard to find the optimal static sparsity pattern for a NN module to minimize the approximation error to the dense model. Second, for each sparsity pattern, we need to take into account hardware block-oriented efficiency (accessing each element in memory takes the same time as accessing the block of adjacent elements~\citep{cook2012cuda}, illustrated in~\ref{fig:hardware}). Common theoretical measures of efficiency (e.g., number of non-zeros, FLOPs) do not map well to modern hardware designed for block computation. Last, every different NN module might require different sparsity patterns, which makes the problem even more complicated.

In our early exploration, we empirically study many sparsity patterns proposed in the literature to find those patterns that can closely approximate the dense model (Details in~\ref{sec:appx_ntk_algorithm}). We found that one sparsity pattern, namely butterfly + low-rank, consistently outperforms the others. This sparsity pattern closely connects to two lines of work in matrix structures: (i) sparse + low-rank matrices, which can capture global and local information~\citep{candes2011robust,udell2019big, scatterbrain}, and (ii) butterfly matrices~\citep{parker1995random,dao2019learning} whose products can tightly represent any sparse matrix~\citep{desa2018two, dao2020kaleidoscope}. Using the fixed sparsity pattern from butterfly matrices, with the addition of a low-rank term, would address two of the three challenges above and yield a simple way to sparsify most NN layers (that are based on matrix multiply). 

However, butterfly matrices are inefficient on modern hardware: (i) they are difficult to parallelize as they contain sequential products of many factors, and (ii) they are not hardware-friendly because the sparsity patterns are not block-aligned. We propose two simple changes to make Butterfly efficient while retaining their favorable properties. Our proposal, Pixelated Butterfly (Pixelfly), combines flat block butterfly and low-rank matrices to yield a simple and efficient sparse training method.

\begin{itemize}[leftmargin=*,nosep,nolistsep]
\item We design an extremely simple sparsity pattern inspired by butterfly + low-rank matrices, which takes into account the hardware's block-oriented efficiency. We propose block butterfly matrices that are efficient as their sparsity patterns align with hardware blocks. We then introduce flat butterfly, a first-order approximation of butterfly with residual connection, that turns the original product of factors into a sum. Flat butterfly matrix multiplications are easy to parallelize. Pixelfly, uses the fixed sparsity pattern from flat \& block butterfly, along with a low-rank term, to produce a sparse network.
\item We prove that block butterfly retains the expressiveness of butterfly matrices and can thus tightly capture sparse matrices. We show that flat butterfly matrices can closely approximate large classes of matrices that butterfly matrices capture. Moreover, we demonstrate that flat block butterfly + low-rank matrices are strictly more expressive than sparse or low-rank matrices alone. Finally, leveraging the recent advance in the neural tangent kernel (NTK), we adapt existing techniques to prove the global convergence of gradient descent on training sparse and wide ReLU networks.
\item Our proposed Pixelfly can be applied to all network modules that rely on matrix multiplication (e.g., linear layer, attention, MLP). To sparsify a full network, we simply need to allocate compute budget for each layer based on matrix and hardware block size.
\end{itemize}

We empirically validate that Pixelfly can speed up the training of models (Transformers, ViT, MLP-Mixer) without quality drop compared to baselines on a wide range of domains and tasks. On CIFAR10/100 \& ImageNet classification, Pixelfly achieve 2.3$\times$ training time speedup compared to dense ViT, MLP-Mixer models, and other sparse training baselines, while preserving the same accuracy. On the WikiText-103 language modeling task, we speed up GPT-2 Medium training by 2.5$\times$ and achieve the same perplexity. On the Long Range Arena benchmark, we maintain the same accuracy as Transformer with $5.2\times$ faster training than a dense model, $2\times$ faster than Sparse transformer, and $6\times$ faster than non-block-aligned sparse methods (Reformer). Our ablation studies highlight the importance of each of our components: our butterfly sparsity improves on existing hand-crafted patterns by up to 2\% of accuracy on ImageNet, our hardware-aware block-sparsity yields up to 5$\times$ speedup, and the balanced compute budget allocation brings 2$\times$ speedup compared to baselines that only sparsify attention.\footnote{Pixelfly code is available at \url{https://github.com/HazyResearch/pixelfly}}

\section{Problem Setting}
\label{sec:problem_formulation}
We first define the problem as sparse matrix approximation with a simple hardware cost model. Then we briefly introduce butterfly and sparse + low-rank matrices. 

\textbf{Problem Formulation:} We focus on the training of GEMM-based models, which can be viewed as a series of matrix multiplies (Given $A, B \in R^{n \times d}$, compute $C = AB^T$). Speeding up training while maintaining model quality can be mapped to finding an approximation procedure $f$ which reduces the time $T$ of computing $C$ while minimizing error $\mathbf{E}[\Vert f(A, B)-AB^T\Vert^2_F]$. Since the hardware is a block device, accessing any individual element within a block of memory is the same as accessing the full block~\citep{cook2012cuda} (\ref{fig:hardware}). A simple cost model of $T$ on hardware with block size $b$ would depend on the number of $b$-blocks being accessed and compute time (formal definition in~\ref{app:problem_formulation}).
Our experiment (\ref{app:study}) reveals that when the non-zeros are grouped into blocks, picking the smallest block size supported by hardware can speed up operations by 10$\times$ compared to sparsity patterns that are not block-aligned.

\textbf{Butterfly, Sparse + Low-rank Matrices:} Butterfly matrices have been used in numerical linear algebra~\citep{parker1995random, li2015butterfly} and machine learning~\citep{mathieu2014fast, jing2017tunable, munkhoeva2018quadrature, dao2019learning, choromanski2019unifying}. They encode the recursive divide-and-conquer structure of the fast Fourier transform (FFT) algorithm~\citep{cooley1965algorithm} and provably capture any sparse matrix with near-optimal space and time complexity. Sparse and Low-rank structures have been studied in Robust PCA~\citep{candes2011robust}, graph clustering~\citep{jalali2011clustering}, and co-variance estimation~\citep{luo2011high}. Recently it has been adopted in attention approximation for Transformers~\citep{scatterbrain}. 

\section{Butterfly matrices and Pixelated Butterfly}
\label{sec:butterfly}

Butterfly matrices~\citep{parker1995random, dao2019learning} are expressive and theoretically efficient. As they contain the set of sparse matrices, we choose to search for the sparsity pattern in this larger class due to their fixed sparsity structure. However, there are three technical challenges. We highlight them here along with our approaches to address them:

\begin{enumerate}[leftmargin=*,nosep,nolistsep]
  \item Slow speed: butterfly matrices are not friendly to modern hardware as their sparsity patterns are not block-aligned, thus are slow. We introduce a variant of butterfly matrices, \emph{block butterfly}, which operate at the block level, yielding a block-aligned sparsity pattern.
  \item Difficulty of parallelization: the sequential nature of butterfly matrices as products of many factors makes it hard to parallelize the multiplication. We propose another class of matrices, \emph{flat butterfly} matrices, that are the first-order approximation of butterfly with residual connections. Flat butterfly turns the product of factors into a sum, facilitating parallelization.
  \item Reduced expressiveness of flat butterfly: even though flat butterfly matrices can approximate butterfly matrices with residual connections, they are necessarily high-rank and cannot represent low-rank matrices~\citep{udell2019big}. We propose to add a low-rank matrix (that is also block-aligned) to flat butterfly to increase their expressiveness.
\end{enumerate}

Combining these three approaches (flat \& block butterfly + low-rank), our proposal (Pixelated Butterfly) is a very simple method to train sparse networks.

\subsection{Block Butterfly Matrices}
\label{sec:block_butterfly}

We propose a block version of butterfly matrices, which is more hardware-friendly than the regular butterfly. The regular butterfly matrices~\citet{dao2019learning, dao2020kaleidoscope} will be a special case of block butterfly with block size $b = 1$. We omit $b$ in the notation if $b = 1$.

\begin{definition} \label{def:bfactor}
  A \textbf{block butterfly factor} (denoted as $\mathbf{B}_{k, b}$) of size $kb$ (where $k \ge 2$) and block size $b$ is a matrix of the form
    \(
        \mathbf{B}_{k, b} = \begin{bmatrix}
            \mathbf{D}_1 & \mathbf{D}_2 \\ \mathbf{D}_3 & \mathbf{D}_4
        \end{bmatrix}
    \)
    where each $\mathbf{D}_i$ is a $\frac{k}{2} \times \frac{k}{2}$ block diagonal matrix of block size $b$ of the form
    $\mathrm{diag} \left( D_{i, 1}, \dots, D_{i, k/2} \right)$ where $D_{i, j} \in \mathbb{R}^{b \times b}$. We restrict $k$ to be a power of 2.
\end{definition}

\begin{definition}
  A \textbf{block butterfly factor matrix} (denoted as $\mathbf{B}_{k}^{(n, b)}$) of size $nb$ with stride $k$ and block size $b$ is a block diagonal matrix of $\frac{n}{k}$ (possibly different) butterfly factors of size $kb$ and block size $b$:
     \[
        \mathbf{B}_{k}^{(n, b)} = \mathrm{diag} \left( \left[ \mathbf{B}_{k, b} \right]_1, \left[ \mathbf{B}_{k, b} \right]_2, \hdots, \left[ \mathbf{B}_{k, b} \right]_\frac{n}{k} \right)
     \]
\end{definition}

\begin{definition} \label{def:bmatrix}
    A \textbf{block butterfly matrix} of size $nb$ with block size $b$ (denoted as $\mathbf{B}^{(n, b)}$) is a matrix that can be expressed as a product of butterfly factor matrices:
    \(
        \mathbf{B}^{(n, b)} = \mathbf{B}_n^{(n, b)} \mathbf{B}_{\frac{n}{2}}^{(n, b)} \hdots \mathbf{B}_2^{(n, b)}.
    \)
    Define $\mathcal{B}_b$ as the set of all matrices that can be expressed in the form $\mathbf{B}^{(n, b)}$ (for some $n$).
\end{definition}

\subsection{Flat butterfly matrices}
\label{sec:flat_butterfly}
In most applications of butterfly matrices to neural networks, one multiplies the $O(\log n)$ butterfly factors. However, this operation is hard to be efficiently implemented on parallel hardware (e.g., GPUs) due to the sequential nature of the operation\footnote{Even with a very specialized CUDA kernel, butterfly matrix multiply ($O(n \log n)$ complexity) is only faster than dense matrix multiply ($O(n^2)$ complexity) for large values of $n$ (around 1024)~\citep{dao2019learning}.}. We instead propose to use a sum of butterfly factors that can approximate the products of the factors. This sum of factors results in one sparse matrix with a fixed sparsity pattern, which yields up to 3$\times$ faster multiplication on GPUs (\ref{sec:appx_benchmark}).

Residual connections have been proposed to connect the butterfly factors~\citep{vahid2020butterfly}. We show that residual products of butterfly matrices have a first-order approximation as a sparse matrix with a fixed sparsity. Let $M$ be a matrix in the set of butterfly matrices $\mathcal{B}$.
In residual form, for some $\lambda \in \mathbb{R}$:

\begin{equation}
  \label{eq:residual_butterfly}
  M = (I + \lambda \mathbf{B}_n^{(n)}) (I + \lambda \mathbf{B}_{n/2}^{(n)}) \dots (I + \lambda \mathbf{B}_2^{(n)}).
\end{equation}

Note that this form can represent the same matrices in the class of butterfly matrices $\mathbf{B}$, since any $\mathbf{B}_k^{(n)}$ contains the identity matrix $I$.

Assuming that $\lambda$ is small, we can expand the residual and collect the
terms\footnote{We make the approximation rigorous in \ref{sec:analysis}.}:

\begin{equation*}
  M = I + \lambda (\mathbf{B}_2^{(n)} + \mathbf{B}_{4}^{(n)} + \dots + \mathbf{B}_n^{(n)}) + \widetilde{O}(\lambda^2).
\end{equation*}

\begin{definition}
  \label{def:flat_butterfly}
  \emph{Flat butterfly} matrices of maximum stride $k$ (for $k$ a power of 2) are those of the form $I + \lambda (\mathbf{B}_2^{(n)} + \mathbf{B}_{4}^{(n)} + \dots + \mathbf{B}_k^{(n)})$.
\end{definition}

Flat butterfly matrices of maximum stride $n$ are the first-order approximation of butterfly matrices in residual form (\ref{eq:residual_butterfly}).
Notice that flat butterfly of maximum stride $k$ are sparse matrices with $O(n \log k)$ nonzeros with a fixed sparsity pattern, as illustrated in~\ref{fig:tradeoff}.
We call this sparsity pattern the \emph{flat butterfly} pattern.

\emph{Flat block butterfly} matrices are block versions of flat butterfly in~\ref{sec:flat_butterfly} (shown in~\ref{fig:tradeoff}).
We empirically validate that flat block butterfly matrices are up to 3$\times$
faster than block butterfly or regular butterfly (\ref{sec:appx_benchmark}).

Since flat butterfly matrices approximate the residual form of butterfly matrices, they have high rank if $\lambda$ is small (\ref{sec:analysis}).
This is one of the motivations for the addition of the low-rank term in our method.

\subsection{Pixelated Butterfly: Flat Block Butterfly + Low-rank for Efficient Sparse Training}
\label{sec:method}

We present Pixelated Butterfly, an efficient sparse model with a simple and fixed sparsity pattern based on butterfly and low-rank matrices. Our method targets GEMM-based neural networks, which are networks whose computation is dominated by general matrix multiplies (GEMM), such as Transformer and MLP-Mixer. As a result, we can view the network as a series of matrix multiplies.

Given a model schema (layer type, number of layers, matrix dimension) and a compute budget, Pixelated Butterfly has three steps: compute budget allocation per layer, sparsity mask selection from the flat butterfly pattern, and model sparsification. We describe these steps in more details:

\begin{enumerate}[leftmargin=*,nosep,nolistsep]
  \item \textbf{Compute budget allocation}: based on our cost model
  (\ref{app:problem_formulation}), given the layer type, number of layers, and matrix dimension, we can find the density (fraction of nonzero weights) of each layer type to minimize the projected compute cost. Continuing our goal for a simple method, we propose to use a simple rule of thumb: allocate sparsity compute budget proportional to the compute fraction of the layer. For example, if the MLP layer and attention layers are projected to takes 60\%
  and 40\% the compute time respectively, then allocate 60\% of the sparsity compute budget to MLP and 40\% to attention. We verify in \ref{sec:appx_method_details} that this simple rule of thumb produces similar results to solving for the density from the cost model.

  \item \textbf{Sparsity mask selection}: given a layer and a sparsity compute budget for that layer, we use one-quarter to one-third of the budget for the low-rank part as a simple rule of thumb. We pick the rank as a multiple of the smallest supported block size of the device (e.g., 32) so that the low-rank matrices are also block-aligned. The remaining compute budget is used to select the sparsity mask from the flat block butterfly sparsity pattern: we choose the butterfly block size as the smallest supported block size of the device (e.g., 32), and pick the maximum stride of the flat block butterfly (\ref{def:flat_butterfly}) to fill up the budget.

  \item \textbf{Model sparsification}: The resulting sparse model is simply a model whose weights or attention follow the fixed sparsity mask chosen in step 2, with the additional low-rank terms (rank also chosen in step 2). In particular, we parameterize each weight matrix\footnote{We describe how to add sparse and low-rank for attention in \ref{sec:appx_method_details}} as: $W = \gamma B + (1 - \gamma) U V^\top$, where $B$ is a flat block butterfly matrix (which is sparse), $U V^\top$ is the low-rank component, and $\gamma$
  is a learnable parameter. We train the model from scratch as usual.
\end{enumerate}

Our method is very simple, but competitive with more complicated procedures that search for the sparsity pattern (\ref{sec:appx_ntk_algorithm}).
We expect more sophisticated techniques (dynamic sparsity, a better approximation of butterfly) to improve the accuracy of the method.

\section{Theoretical analysis}
\label{sec:analysis}

We characterize the expressiveness of the matrices used in our method. In particular, we prove that block butterfly retains the expressiveness of butterfly, and that flat butterfly can accurately approximate the residual form of butterfly. Moreover, flat block butterfly + low-rank (an instance of sparse + low-rank) is more expressive than sparse or low-rank matrices alone. Finally, we analyze the training convergence and generalization of networks with sparse weights. All proofs are in the Appendix.

\subsection{Expressiveness of Block Butterfly} We first prove the expressiveness of block butterfly matrices.

\begin{theorem}
  \label{thm:block_butterfly}
  The set $\mathbf{B}_{2b}$ of $n \times n$ block butterfly matrices with block size $2b$ contains the set $\mathbf{B}_b$ of $n \times n$ block butterfly matrices of block size $b$.
\end{theorem}

By a recursive argument, the set of block butterfly matrices whose block size is a power of 2 contains the set of regular butterfly matrices.

\citet{dao2020kaleidoscope} show that butterfly matrices can tightly represent all structured matrices, such as sparse matrices and many fast transforms. As a result, block butterfly matrices can also represent those structured matrices. In particular,

\begin{corollary}
  \label{cor:block_butterfly_contains_sparse}
  For any constant block size $b$ that is a power of 2, any $nb \times nb$ spare matrix with $s$ nonzeros can be written as products of block butterfly matrices with block size $b$ and their transposes, with $O(s \log n)$ parameters.
\end{corollary}

\subsection{Expressiveness of Flat Butterfly}

We now characterize how the flat butterfly matrices approximate butterfly matrices. In particular, assuming that each butterfly factor has bounded norm, we show that flat-butterfly matrices can accurately approximate the residual form of butterfly with error scaling as $\widetilde{O}(\lambda^2)$.

\begin{theorem}
  \label{thm:flat_butterfly_approx}
  Let $M$ be a matrix of the form in \ref{def:flat_butterfly} where $k = n$, with
  $B_\mathrm{max} := \max_{i} \left\|{\mathbf{B}_i^{(n)}}\right\|_F$ and
  $\left\lvert\lambda\right\rvert \leq \frac{c \sqrt{\epsilon}}{\log n B_\mathrm{max}}$ for some constant $0 < c \leq \frac{1}{2}$ and some $\epsilon > 0$. Then
  \begin{equation*}
    \left\|{M - \left( I + \lambda (\mathbf{B}_2^{(n)} + \mathbf{B}_{4}^{(n)} + \dots + \mathbf{B}_n^{(n)}) \right)}\right\|_F \leq \epsilon.
  \end{equation*}
\end{theorem}

We show that flat butterfly matrices must have high-rank if $\lambda$ is small. This is the motivation for the addition of the low-rank term in Pixelfly (\ref{sec:butterfly}).

\begin{theorem}
  \label{thm:flat_butterfly_rank}
  Let $M$ be as in \ref{eq:residual_butterfly}, with
  $B_\mathrm{max} := \max_{i} \left\|{\mathbf{B}_i^{(n)}}\right\|_F$ and
  $\left\lvert\lambda\right\rvert \leq \frac{c \sqrt{\epsilon}}{\log n B_\mathrm{max}}$ for some constant $0 < c \leq \frac{1}{4}$ and some $\epsilon > 0$. Let $B^\infty_\mathrm{max} = \max_i \left\|{\mathbf{B}_i}\right\|_\infty$.
  Assuming $B^\infty_\mathrm{max} \leq B_\mathrm{max}$.
  Then
  \begin{equation*}
    \mathrm{rank}(I + \lambda (\mathbf{B}_2^{(n)} + \dots + \mathbf{B}_n^{(n)})) = \Omega \left( \left( \frac{B_\mathrm{max}}{B^\infty_\mathrm{max}} \right)^2 \cdot \frac{\log n}{\epsilon \log \left( \frac{B_\mathrm{max}}{B^\infty_\mathrm{max}} \right)} \right).
  \end{equation*}
\end{theorem}

\subsection{Expressiveness of Flat Block Butterfly + Low-rank}

\citet{scatterbrain} prove that there is a natural class of input sequences (generated by a clustering process) whose attention matrix can only be approximated well by sparse + low-rank matrices, and not sparse or low-rank matrices alone. We adapt their technique to show a similar result for the class of matrices we use in Pixelfly.

We require an extra assumption on the clustering process compared
to~\citet{scatterbrain}: the elements in the input sequence form clusters with the same size. Then their attention matrix will have a large block diagonal component well-approximated by flat butterfly, while the rest of the attention matrix is of medium size and is well-approximated by low-rank.

\begin{theorem}[Informal]
  There exists a class of input sequences whose attention matrices are well-approximated by flat block butterfly + low-rank (a special case of sparse + low-rank) but not by sparse or low-rank alone.
\end{theorem}

The formal theorem statement and proof are in \ref{subsec:flat_butterfly_lr_proofs}.

\subsection{Convergence and Generalization of Sparse Networks}

There are natural questions about the training and generalization of sparse models: do they train similarly to dense models, is their generalization close to that of dense models, and can one successfully train them with gradient descent? Our analysis theoretically shows that the answers are yes.

Our analysis relies on the neural tangent kernel (NTK)~\citep{jacot2018neural} of the network. The NTK of two data points $x$ and $y$ measures the similarity between the gradient of the network when evaluated at $x$ compared to the gradient when evaluated at $y$. This kernel governs the dynamics of the neural network output function $f(\cdot, \theta)$ throughout the training and its generalization. We build on the great literature of NTK~\citep{ll18,dzps19,als19_dnn}. The standard result~\citep{sy19} implies the following, if the NTK of the sparse model is close to the NTK of the dense model, then (i) their training convergence speed is similar, (ii) their generalization bounds are similar. For completeness, we state the formal result in \ref{sec:appx_ntk}.

Though this result does not capture the possible regularization effect of sparsity, it shows that sparse models with small NTK difference from dense NTK preserve the generalization ability of dense models, a subject that has been studied more extensively, both from empirical and from theoretical perspectives. We also show that training wide and sparse networks with gradient descent converges globally, similar to the result for wide dense networks~\citep{dzps19,
  als19_dnn} in \ref{sec:gradient_flow}.

\section{Experiments}
\label{sec:experiments}

In this section, our goal is to demonstrate that an extremely simple fixed sparsity pattern can actually speed up sparse model training in wall-clock time without degrading model quality. Specifically, we empirically validate three claims that suggest Pixelfly can improve training speed of different model architectures while retaining model quality on a wide range of domains and tasks.

\begin{enumerate}[leftmargin=*,nosep,nolistsep]
  \item Section~\ref{exp:image}: for image classification tasks, we first show the empirical NTK of flat block butterfly + low-rank sparsity pattern is closer to dense NKT than other baselines. Then we demonstrate our superior end-to-end performance. Specifically, we achieve training speed up on both MLP-Mixer and ViT models by up to 2.3$\times$ wall-clock time with no drop in accuracy compared to the dense model and up to 4$\times$ compared to RigL, BigBird and other sparse baselines.
  \item Section~\ref{exp:lm}: for language modeling and text classification tasks, we can speed up GPT-2 small dense model training by 2.1$\times$, achieving a perplexity of 22.5 on wikitext-103. In addition, on Long Range Arena (LRA) benchmark, we maintain the same accuracy but have $5.2\times$ speed-up in training.
  \item Section~\ref{exp:ablation}: we show the necessity of block flat butterfly and low-rank structures, hardware-alignment and wide coverage of most network layers with ablation studies on these three components of Pixelfly.
 \end{enumerate}

\subsection{Image Classification}
\label{exp:image}

We evaluate the quality and efficiency of Pixelfly through three metrics: (1) distance to training dynamic of the dense model: compare the distance between empirical NTK kernel\footnote{There is an emerging consensus that the NTK is an informative measure of how training and convergence behaviors of two models are similar.} of the models with candidate patterns, including BigBird~\citep{zaheer2020bigbird}, Butterfly~\citep{dao2020kaleidoscope}, and that of the dense model, (2) upstream accuracy: compare the accuracy and training time of the Pixelfly, the dense counterpart, and other baselines on same image classification tasks, (3) downstream accuracy: compare the accuracy of our pretrained Pixelfly and dense model fine-tuned on downstream tasks (\ref{app:throughput}). The empirical NTK of the model with flat block butterfly + low-rank, picked by Pixelfly, is closer to the NTK of the dense model. Pixelfly MLP-mixer and ViT models also retain the same top-1 accuracy of the original dense models while achieving up to 2.3$\times$ speed up.

\textbf{Setup:} We use three popular vision benchmarks, CIFAR-10/100~\citep{krizhevsky2009learning} and ImageNet~\citep{deng2009imagenet}. We choose recent popular Vision Transformer~\citep{dosovitskiy2020image}, T2T-ViT~\citep{yuan2021tokens} and MLP-Mixer~\citep{tolstikhin2021mlp} as representative base models. Their major computation bottlenecks are in different components, e.g. MLP only, attention, or both so we can evaluate the end-to-end applicability of Pixelfly more clearly.

\textbf{Empirical NTK:} To characterize the training dynamic of the sparse networks, we compute the empirical NTK kernels for dense Vision Transformer on CIFAR-100. Then, we show the relative differences between kernels of models with different sparsity patterns and that of the dense one in~\ref{fig:ntk}. Specifically, we pick a popular sparsity pattern combination -- Bigbird pattern~\citep{zaheer2020bigbird} for attention layer and random (magnitude-based sparsity at initialization equals to random) for MLP layer, as a representative baseline. The plot indicates that our designed pattern, flat block butterfly + low-rank is the closest one to that of the dense one among all the patterns. Hence, we expect them to enjoy the most benefits of their dense overparameterized counterparts in real tasks. More details on measuring empirical NTK are covered in the~\ref{app:ntk_exp}.

\textbf{Training from scratch:} We validate that Pixelfly trains up to 2.3$\times$ and 2.0$\times$ faster than dense MLP-Mixer and ViT models from scratch, with the same accuracy under the same setting (batch size, epochs). Specifically, we sparsify the models with Pixelfly and train them on three commonly used vision benchmarking datasets, CIFAR-10/100 and ImageNet. We measure their Top-1 accuracy wall-clock training time. To summarize the general trend,~\ref{table:imagenet} highlights that our sparse vision models consistently retain the accuracy of their dense counterparts in terms of accuracy and achieve training-time speed-up.

Furthermore, we have discussed in~\ref{sec:intro} that current sparse training algorithms aim to dynamic search what could be good sparsity for efficient inference but do not speed up training in wall-clock time. But we still present the comparison results in~\ref{fig:rigl} for completeness. For a fair comparison, we conduct the experiment on Mixer-S/32 model for 100 epochs because RigL aims for sparsity on weights, while we aim for both weights \& attention. As expected, RigL does not speed up training (the pioneering work has unstructured sparsity and does not achieve speed up on GPU) but surprisingly Pixelfly outperforms both dense and RigL in terms of accuracy while achieving $2.1\times$ speedup.

Finally, we compare Pixelfly with BigBird and Sparse Transformer pattern. For a fair comparison, we choose T2T-ViT as the base model because its major bottleneck is on the T2T attention module (our baselines are efficient attention variants). We can see from~\ref{fig:bigbird} that Pixelfly is the only one that can maintain the accuracy and have actual speed up. Further more, Pixelfly speeds up T2T module (large attention) by 1.4$\times$ compare to dense.

\subsection{Language Modeling and Text Classification}
\label{exp:lm}

In this section, we aim to evaluate the effectiveness of Pixelfly in the text domain, on a language modeling task and Long Range Arena (LRA~\citep{tay2020long}) benchmarks. On WikiText-103~\citep{merity2016pointer}, Pixelfly achieves 22.5 perplexity, which is around the same perplexity as GPT-2 small~\citep{radford2019language} but trains 2.1$\times$ faster. On LRA, Pixelfly obtains almost the same accuracy as the full model but gains up to $5.2\times$ speed-up.

\textbf{Setup:} We use WikiText-103 for language modeling and LRA for classification tasks. We use GPT-2 small and vanilla Transformer as the base dense models. The computational bottleneck of GPT-2 small for moderate sequence length, e.g. 512, would be on both attention and MLP layers, while the bottleneck of transformer on LRA task is on attention since the benchmark is designed to evaluate models under long-context scenarios. 

\textbf{GPT-2-Small, Medium on WikiText-103:} We show training GPT-2-Small, Medium and its Pixelfly model from scratch on a commonly used NLP benchmarking dataset, wikiText-103. We measure their perplexity on that dataset, and our training speed up. All setup and finetuning hyperparameters follow the ones in the original paper~\citep{radford2019language}. We present the results in~\ref{table:gpt}. It is not hard to see that Pixelfly models have great advantages in accuracy-efficiency tradeoffs since it maintains the same perplexity as the dense model but achieve up to 2.5$\times$ speed-up in training.

\textbf{Vanilla Transformer on LRA:} We compare vanilla transformer and its Pixelfly models trained from scratch on LRA benchmark. We measure the accuracy, throughput, and training time of both models. Each task has a different sequence length varying between 1024 and 4096. We follow the implementation and experimental setting in~\citep{xiong2021nystr}. We compare the performance of Pixelfly against the dense transformer and report the results in~\ref{table:lra}. We also include the numbers of other baselines from the same repository in the appendix. We can see Pixelfly cause almost no drop in accuracy while achieving 5.2$\times$ speed-up in time.

\subsection{Ablation Study}
\label{exp:ablation}
We conduct ablation studies on each component of Pixelfly (Details in~\ref{app:study}). Specifically, we present (i) how flat block butterfly and low-rank affect the model quality, (ii) how different block size would affect the training speed, (iii) how budget allocation affects the end-to-end speed up.

\textbf{Necessity of Flat Block Butterfly and Low-rank:} (i) We apply different parameter allocation of flat block butterfly and Low-rank component in Pixelfly Mixer-S model on CIFAR-10 under the different density varying in [0.05, 0.1, 0.2]. We found that similar to what was reported in~\citep{scatterbrain}, using around $\frac{1}{4}$ budget on Low-rank and $\frac{3}{4}$ on flat block butterfly achieves the best accuracy. (ii) We also compare Pixelfly with baseline sparsity patterns and show it is $2.7\times$ faster than dense, $3\times$ faster than Butterfly, $3.2\times$ faster than BigBird under $10\%$ density.

\textbf{Block Size:} We study the accuracy-efficiency trade-off for flat block butterfly and random sparsity pattern with different block sizes from 1-32 (~\ref{table:throughput}). We found that first, under the same density, the same sparsity patterns covered with different block sizes could have a big difference in efficiency. Under the same block, the pattern with more locality can be more efficient. Last, the density can seem very small, but actually memory access could be up to $100\%$ of the matrix. Therefore, we always want to make full utilization of the smallest block size that the hardware (or compiler) supported.   

\textbf{Budget Allocation:} We sparsify different components of ViT-small separately, including attention and MLP. We show that their compute ratio is approximately $1:2$ , so if only sparsify one of them, the other one will be the bottleneck preventing end-to-end speed up. Therefore, it is necessary to have an algorithm that can sparsify all layers. 

\section{Related Work}
\textbf{Lottery Ticket Hypothesis.} Models proposed in our work can be roughly seen as a class of manually constructed lottery tickets. Lottery tickets~\citep{frankle2018lottery} are a set of small sub-networks derived from a larger dense network, which outperforms their parent networks. Many insightful studies~\citep{morcos2019one, orseau2020logarithmic, frankle2019stabilizing, frankle2020linear, malach2020proving, pensia2020optimal} are carried out to analyze these tickets, but it remains difficult to generalize to large models due to training cost. In an attempt, follow-up works~\citep{wang2020picking, tanaka2020pruning} show that one can find tickets without training labels. We draw inspiration from one of them, \citet{liu2020finding}, which uses the NTK to avoid using labels in sparsifying networks. Other recent works use specialized hardware to accelerate sparse training~\citep{Goli_2020_CVPR, raihan2020sparse}.

\textbf{Neural Pruning.} Our work is loosely related to neural network pruning. By iteratively eliminating neurons and connections, pruning has seen great success in compressing complex models. Pioneering work~\citep{han2015deep,han2015learning} shows that pruning can produce significantly smaller and faster models for inference. Subsequent methods~\citep{li2016pruning, NIPS2017_a51fb975, dong2017learning,  sanh2020movement, lagunas2021block, zhu2017prune} improve on the quality of the pruned models. While both our and the pruning methods aim to produce sparse models, we target training efficiency, whereas pruning mostly focuses on inference efficiency at the cost of sacrificing training speed.

\textbf{Overparameterized Models and NTK.} Our analysis for sparse model convergence relies heavily on recent advance in neural tangent kernel (NTK)~\citep{jacot2018neural}. NTK is a tool which has been widely used in analyzing overparameterized models' convergence~\citep{ll18,dzps19,als19_dnn,als19_rnn,sy19}, generalization~\citep{all19}, connection to data separability~\citep{os20}, and cost per iteration~\citep{bpsw21}). Deep Double Descent \citep{nakkiran2019deep, d2020triple} conjectures that the generalization error improves as the parameter count grows. It is not surprising that the community is racing to break the record of the largest parameter counts~\citep{radford2019language, brown2020language, dosovitskiy2020image, tolstikhin2021mlp, zhang2021cpm, naumov2019deep, jumper2021highly}. 

We provide extended related work in~\ref{app:related}.

\section{Conclusion}
\label{sec:conclusion}

In our early exploration of many sparsity patterns with complex training procedures, we found that a simple pattern (butterfly + low-rank) consistently (though not always) performed among the best. This motivated us to propose Pixelated Butterfly, a simple and efficient sparse training method. In our quest for simplicity and efficiency, we have chosen to use fixed sparsity that aligns with modern hardware, which was sufficient to yield wall-clock training time speedup without sacrificing accuracy. We are excited about several future directions. Inspired by the remarkable success of model pruning for inference, it is possible that dynamic block sparse mask could be made efficient yet still accurate. Our flat butterfly is a simple first order approximation of the rich class of butterfly matrices, and there could be more sophisticated approximations that retain more expressiveness. Our method is a first step towards the goal of making sparse models train faster than dense models and make them more accessible to the general machine learning community.

\subsubsection*{Acknowledgments}

We thank Laurel Orr, Xun Huang, Sarah Hooper, Sen Wu, Megan Leszczynski, and Karan Goel for their helpful discussions and feedback on early drafts of the paper.

We gratefully acknowledge the support of NIH under No.\ U54EB020405 (Mobilize), NSF under Nos.\ CCF1763315 (Beyond Sparsity), CCF1563078 (Volume to Velocity), and 1937301 (RTML); ONR under No.\ N000141712266 (Unifying Weak Supervision); ONR N00014-20-1-2480: Understanding and Applying Non-Euclidean Geometry in Machine Learning; N000142012275 (NEPTUNE); the Moore Foundation, NXP, Xilinx, LETI-CEA, Intel, IBM, Microsoft, NEC, Toshiba, TSMC, ARM, Hitachi, BASF, Accenture, Ericsson, Qualcomm, Analog Devices, the Okawa Foundation, American Family Insurance, Google Cloud, Salesforce, Total, the HAI-AWS Cloud Credits for Research program, the Stanford Data Science Initiative (SDSI), and members of the Stanford DAWN project: Facebook, Google, and VMWare. The Mobilize Center is a Biomedical Technology Resource Center, funded by the NIH National Institute of Biomedical Imaging and Bioengineering through Grant P41EB027060. The U.S.\ Government is authorized to reproduce and distribute reprints for Governmental purposes notwithstanding any copyright notation thereon. Any opinions, findings, and conclusions or recommendations expressed in this material are those of the authors and do not necessarily reflect the views, policies, or endorsements, either expressed or implied, of NIH, ONR, or the U.S.\ Government. Atri Rudra’s research is supported by NSF grant CCF-1763481.

\end{document}
